%
% 4_RothsteinTrager.tex -- Abschnitt 4 Resultante anwendung in RT
%
% (c) 2026 Jakob Gierer, OST Ostschweizer Fachhochschule
%
% !TEX root = ../../buch.tex
% !TEX encoding = UTF-8
%
\section{Der Satz von Rothstein-Trager
\label{resultante:section:RT}}
\kopfrechts{Der Satz von Rothstein-Trager}
In der Anlaysis verwendet man für die Berechnung von Stammfunktionen
bei Integralen die Partialbruchzerlegung.
Bei dieser Methode müssen die Nullstellen des Nennerpolynoms bekannt sein.
Das Problem ist also, die Residuen des Integranden zu berechnen, ohne
\index{Residuum}%
den Nenner zu faktorisieren.
Die Mathematiker Rothstein und Trager haben unabhängig voneinander den
\index{Rothstein-Trager}%
Satz \ref{resultante:satz:RT} bewiesen.
Der rationale Teil \(a(z)\) muss teilerfremd zu \(b(z)\) sein, mit \(\deg
(a)<\deg (b)\).
Zusätzlich dazu muss das Polynom \(b(z)\) quadratfrei sein, dann sagt
das Liouville-Prinzip die Form
\index{Liouville-Prinzip}%
\begin{equation*}
    \int\frac{a(z)}{b(z)}\,dz=\sum_{k=1}^{m}c_k\log v_k(z)
    \label{resultante:equation:satzRT}
\end{equation*}
der Stammfunktion voraus. Gesucht sind nun die \(c_k\), für die
\begin{equation}
    a(z)-c_k\cdot b'(z) \quad \text{und} \quad b(z)
    \label{resultante:equation:polynome}
\end{equation}
einen nitchttrivialen grössten gemeinsamen Teiler (Teiler \(\neq1\))
\begin{equation}
    v_k(z)=\ggt\left(a(z)-c_k\cdot b'(z),b(z)\right)
    \label{resultante:equation:vks}
\end{equation}
haben.

\begin{satz}[Rothstein-Trager]
\index{Rothstein-Trager, Satz von}%
\index{Satz!von Rothstein-Trager}%
    Die Nullstellen von \(R(c)\) sind genau die Residuen von
\index{Residuum}%
    \(\frac{a(z)}{b(z)}\).
    \label{resultante:satz:RT} 
\end{satz}
\noindent
Die Nullstellen von \(R(c)=\Res_z(a-c\cdot b',b)\) sind die gesuchten \(c_k\).

Dafür eignet sich die Resultante nach Definition
\ref{resultante:definition:resultante} hervorragend.
Sie liefert als Lösung direkt den \(\ggt\) zweier Polynome mit
dazugehöriger Nullstelle.

\subsection{Anwendung der Resultante
\label{resultante:subsection:anwendung}}
Für die Bestimmung der Nullstellen \(c_k\) und des grössten gemeinsamen
Teilers \(v_k(z)\) wird die Resultante
\begin{equation}
    R(c)
    =
    \Res_z\left(a(z)-c\cdot b'(z),b(z)\right)%=0
    \label{resultante:equation:resz}
\end{equation}
benötigt.
In der Resultante \(R(c)\) wird die Variable \(z\) wie in Beispiel
\ref{resultante:beispiel:elimination} eliminiert.
Gesucht sind dann die Nullstellen $c_k$ mit \(R(c_k)=0\).

\begin{beispiel}
    Die Stammfunktion von
\index{Stammfunktion}%
    \begin{equation*}
        \int\frac{1}{z^3+z}\,dz
    \end{equation*}
    soll mit der Methode von Rothstein-Trager und der Resultante bestimmt
    werden.
    Die Bedingungen für die Rothstein-Trager-Methode sind bereits
    erfüllt mit \(a(z)=1\) und \(b(z)=z^3+z\).
    Es müssen daher die \(c_k\) und die dazugehörigen \(v_k(z)\)
    gefunden werden.
    Für die Resultante braucht es
    \begin{equation}
        \begin{aligned}
            a(z)&=1\\
            b(z)&=z^3+z\\
            b'(z)&=3z^2+1,
        \end{aligned}
        \label{resultante:equation:factoren}
    \end{equation}
    sie sieht wie folgt aus:
    \begin{equation}
        \Res_z\left(-3cz^2+1-c,z^3+z\right)=0.
        \label{resultante:equation:resZgut}
    \end{equation}
    Die daraus folgende Sylvester-Matrix ist:
    \begin{equation}
\def\l#1{\bgroup\color{gray!20}#1\egroup}
        \begin{NiceArray}{*{5}{c}}[cell-space-limits=3pt]
        -3c   & 0     & 1-c & \l{0} & \l{0} \\
        \l{0} & -3c   & 0   & 1-c   & \l{0} \\
        \l{0} & \l{0} & -3c & 0     & 1-c   \\
        1     & 0     & 1   & 0     & \l{0} \\
        \l{0} & 1     & 0   & 1     & 0
        \CodeAfter
            \tikz {
                \draw ([xshift=-3pt,yshift=1pt]1-|1) rectangle ([xshift=4pt,yshift=-1pt]6-|6);
                \draw[dashed] ([xshift=-3pt,yshift=1pt]4-|1) -- ([xshift=4pt,yshift=1pt]4-|6);
            }
        \end{NiceArray}
        \;\to\;
        \begin{NiceArray}{*{5}{c}}[cell-space-limits=3pt]
        \;1\;     & \;0\; & \frac{c-1}{3c} & \l{0}          & \l{0} \\
        \l{0} & 1     & 0              & \frac{c-1}{3c} & \l{0} \\
        \l{0} & \l{0} & 1              & 0              & \frac{c-1}{3c} \\
        1     & 0     & 1              & 0              & \l{0} \\
        \l{0} & 1     & 0              & 1              & 0
        \CodeAfter
            \tikz {
                \draw ([xshift=-4pt,yshift=1pt]1-|1) rectangle ([xshift=3pt,yshift=-1pt]6-|6);
                \draw[dashed] ([xshift=-4pt,yshift=0.5pt]4-|1) -- ([xshift=3pt,yshift=0.5pt]4-|6);
            }
        \end{NiceArray}\,\,\,.
        \label{resultante:equation:sylGutStep1}
    \end{equation}
    Nach Anwendung der Zeilenoperationen erhält man das erste Tableau.
    Auf dem kleineren Bereich kann nochmals eine Zeilenoperation
    durchgeführt werden:
    \begin{equation}
\def\l#1{\bgroup\color{gray!20}#1\egroup}
        \begin{NiceArray}{*{5}{c}}[cell-space-limits=3pt]
        \;1\; & \,\;0\;\, & \frac{c-1}{3c}  & \l{0}           & \l{0} \\
        \l{0} & 1         & 0               & \frac{c-1}{3c}  & \l{0} \\
        \l{0} & \l{0}     & 1               & 0               & \frac{c-1}{3c}\\
        \l{0} & \l{0}     & \frac{2c+1}{3c} & 0               & \l{0} \\
        \l{0} & \l{0}     & \l{0}           & \frac{2c+1}{3c} & 0
        \CodeAfter
            \tikz {
                \draw ([xshift=-4pt,yshift=1pt]1-|1) rectangle ([xshift=3pt,yshift=-1pt]6-|6);
                \draw[dashed] ([xshift=1pt]3-|3) -- ([xshift=1pt,yshift=-1pt]6-|3);
                \draw[dashed] ([xshift=1pt]3-|3) -- ([xshift=3pt]3-|6);
                \draw[dashed] ([xshift=1pt]4-|3) -- ([xshift=3pt]4-|6);
            }
        \end{NiceArray}
        \;\to\;
        \begin{NiceArray}{*{3}{c}}[cell-space-limits=3pt]
        \,1\, & \;0\; & \l{0}         \\
        \l{0} & 1     & 0             \\
        \l{0} & \l{0} & \frac{c-1}{3c}
        \CodeAfter
            \tikz {
                \draw ([xshift=-3pt,yshift=1pt]1-|1) rectangle ([xshift=3pt,yshift=-1pt]4-|4);
                \draw[dashed] ([xshift=-3pt]3-|1) -- ([xshift=3pt]3-|4);
            }
        \end{NiceArray}\,\,\,.
        \label{resultante:equation:sylGutStep2}
    \end{equation}
    Das Schlusstableau hat als entscheidendes Element:
    \begin{equation}
        \begin{aligned}
            \frac{c-1}{3c}&\overset{!}{=}0\\
            c_1&=1,
        \end{aligned}
        \label{resultante:equation:solGut1}
    \end{equation}
    für \(c_1=1\) lässt sich im Schlusstableau aus
    (\ref{resultante:equation:sylGutStep2}) für \(v_1(z)\) der \(\ggt=z\)
    ablesen.
    Findige Leser konnten bereits im vorherigen Tableau aus
    (\ref{resultante:equation:sylGutStep2}) den Algorithmus beenden,
    nämlich mit:
    \begin{equation}
        \begin{aligned}
            \frac{2c+1}{3c}&\overset{!}{=}0\\
            2c+1&=0\\
            2c&=-1\\
            c_2&=-\frac{1}{2},
        \end{aligned}
        \label{resultante:equation:solGut2}
    \end{equation}
    für \(c_2=-\frac{1}{2}\) lässt sich ebenfalls im zweitletzten
    Tableau aus (\ref{resultante:equation:sylGutStep2}) für \(v_2(z)\)
    der \(\ggt=z^2+1\) ablesen.

    Setzt man die gefundenen Nullstellen mit ihrem jeweils zugehörigen
    grössten gemeinsamen Teiler zusammen, erhält man die Stammfunktion:
    \begin{equation*}
        \int\frac{1}{z^3+z}\,dz=c_1\log v_1(z)+c_2\log v_2(z)
        =
        1\log z-\frac{1}{2}\log(z^2+1). \qedhere
    \end{equation*}
    \label{resultante:beispiel:gutesIntegral}
\end{beispiel}

\begin{beispiel}
    Die Methode soll mit einem weiteren Beispiel verdeutlicht werden.
    Das Integral
    \begin{equation*}
        \int\frac{1}{\log x}\,dx
    \end{equation*}
    soll gelöst werden.
    Dafür muss \(F(\vartheta)=\mathbb{Q}(x)\) mit dem logarithmischen
    Element \(\vartheta=\log x\) erweitert werden.
    Der Integrand wird dadurch zu
    \begin{equation*}
        f=\frac{1}{\vartheta}
    \end{equation*}
    mit
    \begin{equation}
        \begin{aligned}
            a(\vartheta)&=1\\
            b(\vartheta)&=\vartheta\\
            Db(\vartheta)&=\frac{1}{x}.
        \end{aligned}
    \end{equation}
    Das \(D\) steht hier für die Derivation nach \(x\).
    Die Resultante daraus sieht wie folgt aus:
    \begin{equation}
        R(c)
        =
        \Res_\vartheta(a(\vartheta)-cDb(\vartheta),b(\vartheta))
        =
        \Res_\vartheta\biggl(1-c\frac{1}{x},\vartheta\biggr)=0,
        \label{resultante:equation:resTbad}
    \end{equation}
    das erste Polynom in der Resultante ist vom Grad \(0\) und das zweite
    vom Grad \(1\).
    Die daraus entstehende \(1\times1\)-Sylvester-Matrix im Tableau ist dann
    \begin{equation}
        \begin{NiceArray}{c}[cell-space-limits=3pt]
        1-c\frac{1}{x} \\
        \CodeAfter
            \tikz \draw ([xshift=-3pt,yshift=2pt]1-|1) rectangle ([xshift=3pt,yshift=-2pt]2-|2);
        \end{NiceArray}
        \,\,\,\,\,\,.
        \label{resultante:equation:sylBad}
    \end{equation}
    Der Algorithmus für dieses Tableau endet, bevor er überhaupt das
    erste Mal angewandt wird.
    Das kritische Element ist:
    \begin{equation}
        \begin{aligned}
            1-\frac{c}{x}&\overset{!}{=}0\\
            \frac{c}{x}&=1\\
            c&=x,
        \end{aligned}
        \label{resultante:equation:solBad}
    \end{equation}
    die Nullstelle \(c\) muss für die Methode von Rothstein-Trager eine
    Konstante sein (\(Dc\overset{!}{=}0\)).
    Hier ist
    \begin{equation*}
        Dc=Dx=1\overset{!}{=}0 \,\text{\lightning}
    \end{equation*}
    Daraus folgt: das Integral ist nicht elementar lösbar.
    Wie Beispiel \ref{buch:intalg:risch:bsp:1/logx} und Definition
    \ref{buch:intalg:risch:def:integrallogarithmus} zeigt, ist das der
    Integrallogarithmus.
    \label{resultante:beispiel:badIntegral}
\end{beispiel}

Sobald die Nullstellen \(c_k\) aus der Resultante keine Konstanten sind,
folgt daraus, dass solche Integrale nicht elementar lösbar sind.
